
\newpage

\Heading{Bibliographie}{0}

{\nspace{
\ref{}
H. Bass.
{\it Generators and relations for cyclotomic
units}, Nagoya Math. J. {\bf 27} (1966),\break
p.~401--407.

\medskip
\ref{}
A. Beilinson, J. Bernstein, et P. Deligne.  
{\it Faisceaux pervers},  in: 
 Analyse et topologie sur
les espaces singuliers, Ast\'erisque 100,
SMF (1982).

\medskip
\ref{}
S. Bloch.
{\it Algebraic cycles and algebraic
$K$-theory},
Adv. in Math. {\bf 61}, 3 (1986), p.~267--304.

\medskip
\ref{}
\hbox to 35pt{\hrulefill}~.
{\it The moving lemma for higher Chow groups},
J. Alg. Geom. {\bf 3}, 3 (1994), p.~537--568.

\medskip
\ref{}
K. T. Chen.
{\it Iterated integrals of differential forms
and loop space homology}, Ann. of Math. {\bf
97} (1973), p.~217--246.

\medskip
\ref{}
\hbox to 35pt{\hrulefill}~.
{\it Reduced Bar Constructions on de Rham
complexes}, in: Algebra, Topology and Category
Theory, a collection of papers in honor of
Samuel Eilenberg, p.~19--32,
Academic Press (1976).


\medskip
\ref{}
P. Deligne.
{\it Le groupe fondamental de la droite
projective moins trois points} in:
Galois group over $\dbQ$, MSRI Publ. 16,
p.~79--313, Springer Verlag (1989).

\medskip
\ref{}
\hbox to 35pt{\hrulefill}~.
{\it Cat\'egories tannakiennes} in:
Grothendieck Festschrift vol. 2, p.~111--195.
Progress in Math. 87, Birkh\"auser (1990).

\medskip
\ref{}
M. Demazure et P. Gabriel.
{\it Groupes alg\'ebriques}, Masson (1970).

\medskip
\ref{}
A.~B.~Goncharov.
{\it Polylogarithms in arithmetic and
geometry}, Proc. ICM Zurich 1994, p.~374--387,
Birkhauser.

\medskip
\ref{}
A.~B.~Goncharov.
{\it The dihedral Lie algebras and Galois
symmetries of $\pi_1(\dbP^1-\{0,\infty\}\CUP
\mu_N)$}, Duke Math. J. {\bf 110} 3 (2001), p.~397--487.

\medskip
\ref{}
R. Hain and M. Matsumoto.
{\it Weighted completion of Galois groups and
Galois actions on the fundamental group of
$\dbP^1-\{0,1,\infty\}$}, preprint (2001).

\medskip
\ref{}
R. Hain and S. Zucker.
{\it Unipotent variations of mixed Hodge
structure}, Inv. Math. {\bf 88}, (1987)
p.~83--124.

\medskip
\ref{}
M. Hanamura.
{\it Mixed motives and algebraic cycles I.},
Math. Res. Lett. {\bf 2}, 6 (1995),
p.~811--821.
See also II, preprint (1996).

\medskip
\ref{}
A. Huber.
{\it Mixed Motives and their Realization in
Derived Categories}, Lecture Notes in Math.
1604, Springer Verlag (1995).

\medskip
\ref{}
\hbox to 35pt{\hrulefill}~.
{\it Realization of Voevodsky's motives}, J.
Alg. Geom. {\bf 9} (2000), p.~755--799.

\medskip
\ref{}
U. Jannsen.
{\it Mixed motives and algebraic $K$-theory},
Lecture Notes in Math. 1400, Springer Verlag
(1990).

\medskip
\ref{}
D. Kubert.
{\it The universal ordinary distribution},
Bull. SMF {\bf 107} (1979), p.~179--202.

\medskip
\ref{}
M. Levine.
{\it Tate motives and the vanishing
conjectures for algebraic $K$-theory}, in:
Algebraic $K$-theory and algebraic topology
(Lake Louise, 1991), p.~167--188.
NATO Adv. Sci. but. Ser. C Math. Phys. {\bf
407}, Kluwer (1993).

\medskip
\ref{}
\hbox to 35pt{\hrulefill}~.
{\it Bloch's higher Chow groups revisited}, in:
$K$-theory (Strasbourg 1992), p.~235--320,
Ast\'erisque 226, SMF (1994).

\medskip
\ref{}
\hbox to 35pt{\hrulefill}~.
{\it Mixed Motives}, Math. Surveys and
Monographs {\bf 57}, AMS (1998).

\medskip
\ref{}
Chr. Reutenauer.
{\it Free Lie algebras}, LMS monographs new ser.
{\bf 7},
Oxford University Press (1993).

\medskip
\ref{}
T. Terasoma.
{\it Multiple zeta values and mixed Tate
motives}.

\medskip
\ref{}
V. Voevodsky.
{\it Triangulated categories of motives over a
field}, in: Cycles, transfer and motivic
homology theories, p.~188--238,
Ann. of Math. Studies 143, Princeton University
Press (2000).

\medskip
\ref{}
L. C. Washington.
{\it Introduction to cyclotomic fields --
graduate texts in Math.} {\bf 81}, Springer
Verlage (1997).


\medskip
\ref{}
Z. Wojtkowiak.
{\it Cosimplicial objects in algebraic
geometry}, in Algebraic $K$-theory and
algebraic topology (Lake Louise, 1991), p.~287--327.
NATO Adv. Sci. but. Ser. C Math. Phys. {\bf
407}, Kluwer (1993).

}}

